



The aim of this appendix is to give an explicit description of the symmetric monoidal structure on the category $\FPMp$ of Poincar\'e form parameters obtained in Proposition \reffour{proposition:sym_mon_structure_on_cat_of_form_parameters}, and of the resulting algebras and modules in this category, in terms of the classical definition of form parameters of Bak~\cite{bak-form-parameter} and  Schlichting~\cite{SchlichtinghigherI}.

\subsection{The box product of Mackey functors}

Recall that the $\infty$-category $\Sp^{\gCt}$ admits a $t$-structure which is pointwise if we identify $\Sp^{\gCt}$ with the category of spectral $\Ct$-Mackey functors, and that the heart of this t-structure is the category of  $\Ct$-Mackey functors valued in abelian groups. The symmetric monoidal structure on $\Sp^{\gCt}_{\geq 0}$ given by Day-convolution then localises onto the heart to a symmetric monoidal structure again given by the Day-convolution of $\Ct$-Mackey functors valued in abelian groups, which is usually called the box-product. 

Recall that a $\Ct$-Mackey functor can be explicitly described as a pair of abelian groups $A$ and $B$, an involution $\sigma$ on $A$, and group homomorphisms
\[
A_{\Ct}\xrightarrow{\tau}B\xrightarrow{\rho}A^{\Ct}
\]
whose composite is the map induced by $\id_A+\sigma$. The map $\rho$ is called the restriction and $\tau$ the transfer. We call $A$ the underlying abelian group and $B$ the fixed-point datum of the Mackey functor. The Mackey functors form the objects of a category that we denote by $\Mack$, where the morphisms are the pairs of group homomorphisms which commute with the involutions, transfers and restrictions.  Under this description of  $\Ct$-Mackey functors, the Day-convolution product takes the following form, see e.g. \cite{Bouc}*{Proposition 1.5.1}.



\begin{definition}
\label{definition:box-mackey}%
The box-product of two $\Ct$-Mackey functors is given by
\[
\big(A_{\Ct}\xrightarrow{\tau}B\xrightarrow{\rho}A^{\Ct}\big)
\boxutens
\big(A'_{\Ct}\xrightarrow{\tau'}B'\xrightarrow{\rho'}(A')^{\Ct}\big)
=\big((A\utens A')_{\Ct}\xrightarrow{}(A\utens A'\oplus B\utens B')/\!\!\sim\to(A\utens A')^{\Ct}\big)
\]
where the quotient is by the subgroup generated by:
\begin{enumerate}
\item $a\otimes a'-\sigma(a)\otimes \sigma'(a')$ for all $a\in A, a'\in A'$,
\item  $b\otimes \tau(a')-\rho(b)\otimes a'$, for all $a'\in A', b\in B$,
\item $\tau(a)\otimes b'-a\otimes \rho(b')$, for all $a\in A, b'\in B'$.
\end{enumerate}
The abelian group $A\utens A'$ has the involution $\sigma\utens\sigma'$,
the restriction map sends $[a\otimes a']$ to $a\otimes a'+\sigma(a)\otimes \sigma'(a')$ and $[b\otimes b']$ to $\rho(b)\otimes \rho(b')$, and the transfer sends $a\otimes a'$ to the equivalence class $[a\otimes a']$.
The unit of the resulting symmetric monoidal structure is the Burnside Mackey functor
\[
\mathbb{A}=\big(\mathbb{Z}\xrightarrow{\tau}\mathbb{A}(\Ct)\xrightarrow{\rho}\mathbb{Z}\big)
\]
where $\mathbb{Z}$ has the trivial involution, and $\mathbb{A}(\Ct)$ is the Burnside ring of $\Ct$ defined as the free abelian group on the (two) isomorphism classes of transitive $\Ct$-sets. The restriction sends such an isomorphism class to its number of elements, and the transfer sends $1$ to the class of the free transitive $\Ct$-set $\Ct$.
\end{definition}








\begin{remark}
\label{remark:green}%
The category of monoids in $\Mack$ with respect to the box-product is equivalent to the category of the so called Green functors, see e.g. \cite{Bouc}*{\S2.3}. A Green functor consists of a Mackey functor as above where $A$ and $B$ are equipped with a ring structure, such that $\sigma$ and $\rho$ are ring homomorphisms, and such that the transfer $\tau$ satisfies the Frobenius reciprocity identity
\[
\tau(\rho(c)\cdot a\cdot \rho(b))=c\cdot \tau(a)\cdot b
\]
for all $a\in A$ and $b,c\in B$. %
Similarly, a (left) module over such a Green functor consists of a Mackey functor 
\[
\big(N_{\Ct}\xrightarrow{\tau}P\xrightarrow{\rho}N^{\Ct}\big)
\]
with an $A$-module structure on $N$ and a $B$-module structure on $P$ such that
\begin{enumerate}
\item
$\sigma(a\cdot n)=\sigma(a)\cdot\sigma(n)$.
\item
$\rho(b\cdot p)=\rho(b)\cdot\rho(p)$
\item
$\tau(\rho(b)\cdot n)=b\cdot\tau(n)$
\item
$\tau(b\cdot \rho(n))=\tau(b)\cdot n$
\end{enumerate}
for all $a\in A$, $b\in B$, $n\in N$ and $p\in P$.
\end{remark}



\begin{example}
\label{example:mackey}%
Given a ring $R$, there is a monoid $\Nmack(R)$ in $\Mack$ defined by
\[
\Nmack(R)=\big((R\utens R)_{\Ct}\xrightarrow{\tau} W(R)\xrightarrow{\rho}(R\utens R)^{\Ct}\big)
\]
where the involution on $R\utens R$ flips the tensor factors, and $W(R)$ is the set $(R\times (R\utens R)_{\Ct})$ with the ring structure defined by
\begin{align*}
(r,x)+(r',x')&=(r+r',x+x'-[r\otimes r'])
\\
(r,[a\otimes b])\cdot(r',[a'\otimes b'])&=(r\cdot r',[r\otimes r\cdot  a'\otimes b']+[a\otimes b\cdot r'\otimes r']+[a\otimes b\cdot a'\otimes b'+a\otimes b\cdot b'\otimes a']).
\end{align*}
The transfer and restriction are defined by $\tau(x)=(0,x)$ and $\rho(r,[a\otimes b])=r\otimes r+a\otimes b+b\otimes a$. In the terminology of  \cites{DKNP1,DKNP2}, $W(R)$ is the ring of ``$2$-truncated $2$-typical Witt vectors of $\ZZ$ with coefficients in $R$''. Moreover, there is an isomorphism of monoids
\[
\Nmack(R)\cong \pi^{\mack}_0NR
\]
in $\Mack$, see \cite{DKNP2}*{Theorem C}.

In fact, $\Nmack$ arises from a symmetric monoidal functor $\Nmack\colon\Ab\to\Mack$. As such, there are canonical isomorphisms of monoids $\Nmack(R\utens R')\cong \Nmack(R)\boxutens \Nmack(R')$ (see \cite{Read}*{Proposition 4.27}) and $\Nmack(\mathbb{Z})\cong\mathbb{A}$, where the last one is induced by the isomorphism
\[
\mathbb{A}(\Ct)\stackrel{\cong}{\longrightarrow} W(\mathbb{Z})
\]
that takes the one-elements $\Ct$-set $1$ to $(1,0)$, and the free transitive  $\Ct$-set $\Ct$ to $(0,1)$.
\end{example}



\subsection{Rings and modules in form parameters}


Recall from Proposition \reffour{proposition:sym_mon_structure_on_cat_of_form_parameters} and its proof that in our framework a form parameter consists of a ring $R$ and a $\Ct$-Mackey functor $M$ which is a module over $\Nmack(R) = \pi^{\mack}_0NR$. %
Under the characterisation of $\pi^{\mack}_0NR$ and that of modules in $\Mack$ from Remark \reffour{remark:green}, one can unravel this definition
to the following, classical definition from \cite{SchlichtinghigherI} (see also Remark \refone{remark:form-parameters-and-modules-with-genuine-involution}).

\begin{definition}
A form parameter over a discrete ring $R$ consists of a $\Ct$-Mackey functor
\[
M_{\Ct}\xrightarrow{\tau}Q\xrightarrow{\rho}M^{\Ct}
\]
together with the structure of an $R\utens R$-module with involution on $(M,\sigma)$, an action on $Q$ by the multiplicative monoid of $R$ through group homomorphisms, such that $\rho$ and $\tau$ are $R$-equivariant and
\[
(r+s)x=rx+sx+\tau((r\otimes s)\rho(x))
\]
for all $r,s\in R$ and $x\in Q$. We denote such a form parameter by $\fpm$, and typically display it by its underlying Mackey functor as in the diagram above.
We say that a form parameter is Poincar\'e if $M$ is projective as an $R$-module and the map
\[
R\longrightarrow \hom_R(M,M)
\]
of Definition \refone{remark:invertible-criterion} is an isomorphism (so that $M$ is invertible).
\end{definition}

Here, the multiplicative action of $R$ on $Q$ corresponds to the action of the elements of the form $(r,0)\in W(R)$ under the $W(R)$-module structure on $Q$, and by the definition of the sum in $W(R)$ we find the relation
\begin{align*}
(r+s,0)\cdot x&=((r,0)+(s,0)+(0,[r\otimes s])))\cdot x=((r,0)+(s,0)+(0,[r\otimes s])))\cdot x
\\&=(r,0)\cdot x+(s,0)\cdot x+\tau([r\otimes s])\cdot x=(r,0)\cdot x+(s,0)\cdot x+\tau(r\otimes s\rho(x)).
\end{align*}







\begin{definition}
We let $\FPM$ be the category of form parameters, whose objects are pairs $(R,\fpm)$ of a ring $R$ and a form parameter $\fpm$ over $R$. A morphism from  $(R,\fpm)$ to  $(R',\fpm')$ consists of a ring homomorphism $f\colon R\to R'$ and a morphism of form parameters $g\colon \lambda\to f^\ast \fpm'$ over $R$, which is defined as a pair of morphisms $M\to M'$ and $Q\to Q'$ which commute with $\sigma$, $\rho$ and $\tau$, and with the respective actions of $R\utens R$ and $R$. 

We let $\FPMp$ be the subcategory of Poincar\'e form parameters, where the morphisms are those such that the underlying map of modules induces an isomorphism  $R'\utens M\cong M'$ (cf. Lemma~\refone{lemma:Poincresscalars}).
\end{definition}




By the discussion at the last part of \S\reffour{subsection:multiplicative-GW}, Poincaré form parameters on $R$ correspond to Poincaré structures $\QF$ on $\Dperf(R)$ such that $\Dual_{\QF}$ preserves $\Proj(R)$ and $\QF$ takes discrete values on $\Proj(R)$. 
The symmetric monoidal structure on $\FPMp$ from Proposition \reffour{proposition:sym_mon_structure_on_cat_of_form_parameters} is by construction that of algebras over $\pi^{\mack}_0N$. We thus obtain the following mild extension of the tensor product of form rings of  \cite{SchlichtinghigherI}*{\S 4} (see also Example \reffour{example:formrings} below).


\begin{proposition}
\label{prop:tensormackeyvstensorfpm}%
 The box product of Mackey functor lifts to a symmetric monoidal structure on the category $\FPM$ of form parameters, whose tensor product sends an $N^{\mack}(R)$-module $X$ and an $N^{\mack}(S)$-module $Y$ in $\Mack$ to $X\boxutens Y$, which we regard as a $N^{\mack}(R\utens S)$-module via the monoidal structure of $N^{\mack}$.

This symmetric monoidal structure restricts to that of Proposition \reffour{proposition:sym_mon_structure_on_cat_of_form_parameters} on the subcategory $\FPMp$ of Poincar\'e form parameters.

\qed
\end{proposition}

We can now conveniently unpack the definition of a monoid and of a module in  $\FPM$ as follows.

\begin{lemma}%
\label{lemma:monoidsFPM}%
A monoid in $\FPM$ consists of a commutative ring $R$, a Green functor $
X=\big(M_{\Ct}\xrightarrow{\tau}Q\xrightarrow{\rho}M^{\Ct}\big)
$  (i.e. a monoid in $\Mack$)
where $M$ has the structure of an $R\utens R$-algebra with involution, and a multiplicative unital map $n\colon R\to Q$ which satisfies 
\begin{enumerate}
\item $n(0)=0$,
\item $\rho n(r)=r\otimes r\cdot 1$
\item $n(r+s)=n(r)+n(s)+\tau(r\otimes s\cdot 1)$
\end{enumerate}
for all $r,s\in R$, where $1\in M$ is the multiplicative unit. Then $X$ is a form parameter with the action $r\cdot q=n(r)q$ for $r\in R$ and $q\in Q$. 

A monoid in $\FPMp$ is a monoid in $\FPM$ whose underlying form parameter is Poincar\'e and such that the canonical map of $R$-modules $R\to M$ induced by the unit of $M$ is an isomorphism.
\end{lemma}

\begin{proof}
By the definition of the tensor product of form parameters, if $X$ is a form parameter over $R$ equipped with the structure of a monoid, then in particular $R$ is a monoid in rings, that is a commutative ring. Then the remaining structure of the monoid is that of an $\Nmack(R)$-algebra, that is the structure of a Green functor on $X$ and a morphism of Green functors $\Nmack(R)\to X$. On underlying rings this map gives $M$ the structure of an $R\utens R$-algebra with involution. It remains to see that the ring homomorphim $g\colon W(R)\to Q$ from the map of Green functors is equivalent to a map $f$ with the properties above. 
Since the elements $(0,x)\in W(R)$ are in the image of the transfer, $g$ is determined by its values on the elements $g(r,0)$. This is the multiplicative map $n\colon R\to Q$ given by $n(r)=g(r,0)$, and the last identity follows from the additivity of $g$, since
\[
n(r+s)=g(r+s,0)=g((r,0)+(s,0)+(0,r\otimes s))=n(r)+n(s)+\tau(r\otimes s\cdot 1)
\]
for all $r,s\in R$. The second identity holds by the compatibility with the restriction maps, and $n$ preserves $0$ since $g$ does.
The condition that $R\cong M$ for monoids in $\FPMp$ is exactly the condition for the unit map $\mathbb{A}\to X$ to lie in $\FPMp$.
\end{proof}




\begin{example}%
\label{example:formrings}%
A commutative form ring in the sense of Schlichting \cite{SchlichtinghigherI}*{Definitions 4.1, 4.6} is exactly a commutative monoid in $\FPMp$, where one defines $n\colon R\to Q$ by $n(r)=r\cdot 1$, where $1$ is the ring unit of $Q$.
\end{example}

\begin{example} 
Recall that a Tambara functor for the group $\Ct$ is a Mackey functor
\[
A_{\Ct}\xrightarrow{\tau}B\xrightarrow{\rho}A^{\Ct}
\]
together with the structure of commutative rings on $A$ and $B$ making this into a Green functor, and a multiplicative unital map $n\colon A\to B$ which satisfies exactly the conditions of Lemma \reffour{lemma:monoidsFPM}. Thus, such a Tambara functor defines a commutative monoid in $\FPM$ with underlying ring $A$.
\end{example}


\begin{lemma}%
\label{lemma:modulesFPM}%
Let $X$ be a monoid in $\FPM$ as in Lemma \reffour{lemma:monoidsFPM}.
A module over $X$ consists of an $R$-algebra $S$ and a module 
$
Y=\big(N_{\Ct}\xrightarrow{\tau}P\xrightarrow{\rho}N^{\Ct}\big)
$
over the Green functor $X$ as in Remark \reffour{remark:green} which is additionally a form parameter over $S$, where $N$ is a module over $M$ in $R\utens R$-modules, and which satisfies the following compatibility condition with the form parameter structure:
\[
(n(r)q)(s\cdot p)=(rs)\cdot (qp)
\]
for all $r\in R, s\in S, q\in Q$ and $p\in P$, where $r\cdot p$ denotes the action of the form parameter $Y$.

If $X$ is a monoid in $\FPMp$, a module in $\FPMp$ over $X$ is a module in $\FPM$ over $X$ whose underlying form parameter is in $\FPMp$. 
\end{lemma}
\begin{proof}
We start by observing the the module action is given by a morphism of form parameters
\[
(R\utens S,X\boxutens Y)=(R,X)\utens (S,Y)\longrightarrow (S,Y).
\]
The map in the first component givens $S$ the structure of an $R$-algebra. The map in the second component is a map $X\boxutens Y\to Y$ of $\Nmack(R\utens S)$-modules, where $Y$ is regarded as a $\Nmack(R\utens S)$-module via the action of $R$ on $S$. This last linearity gives the compatibility of the module structure with the quadratic action, again using that the action of $W(R\otimes S)$ is determined by that of the elements of the form $(r\otimes s,0)$.

If $X$ is a monoid in $\FPMp$, so that we may assume that $M=R$, the condition for the action map of $X$ on $Y$ to lie in $\FPMp$ is that it induces an isomorphism
\[
S\otimes_{R\utens S}(R\utens N)\to N
\]
which is automatic since $N$ is a projective $S$-module. 
\end{proof}


\begin{example}
The tensor unit $\mathbb{A}$ is clearly a module over $\Nmack(\mathbb{Z})$ (it is isomorphic to $\Nmack(\mathbb{Z})$ by the monoidality of $\Nmack$). The multiplicative action of $\mathbb{Z}$ on $\mathbb{A}(\Ct)$ equipping it with the structure of $\Nmack(\mathbb{Z})$-module is given by
\[
a\cdot x=(a+\frac{a^2-a}{2}\Ct) x
\]
where $x$ is the isomorphism class of a $\Ct$-set and $a\in\mathbb{Z}$. This formula comes from the orbit decomposition of the set $k\times k$ with $k^2$-elements, $k$ a natural number, equipped with the $\Ct$-action that flips the product factors.

Since $\mathbb{A}$ is the tensor unit it is also a monoid in $\FPMp$. The ring structure on $\mathbb{A}(\Ct)$ is the usual one of the Burnside ring defined from the cartesian product of $\Ct$-sets, and the map $n\colon \mathbb{Z}\to \mathbb{A}(\Ct)$ is
\[
n(a)=a+\frac{a^2-a}{2}\Ct.
\]
Moreove, again because $\mathbb{A}$ is the tensor unit, any form parameter is uniquely a module over $\mathbb{A}$.
\end{example}


\begin{example}
\label{example:fixeptmonoid}%
Let $R$ be a ring and $M$ an $R\utens R$-module with involution $\sigma$. The fixed point Mackey functor
\[
M^{\sym}=\big(M_{\Ct}\xrightarrow{\id+\sigma}M^{\Ct}\xrightarrow{\id}M^{\Ct}\big)
\]
is a form parameter over $R$, called the \defi{symmetric} form parameter, where the multiplicative action of $R$ on $M^{\Ct}$ is defined by $r\cdot m=(r\otimes r) m$, for all $r\in R$ and $m\in M^{\Ct}$. The corresponding Poincaré structure on $\Dperf(R)$ is the genuine symmetric structure.

If $R$ is moreover commutative and $M$ is a commutative $R\utens R$-algebra with involution, then $M^{\sym}$ is a commutative monoid in $\FPM$ where the multiplication on $M^{\Ct}$ is restricted from that of $M$, and the map $n\colon R\to M^{\Ct}$ is defined by
\[
n(r)=r\otimes r\cdot 1,
\]
where $1$ is the unit of $M$. In particular if $R$ is a commutative ring with involution then $R^{\sym}$ is a commutative monoid in $\FPMp$.
\end{example}



\begin{example}
Let $S$ be a ring and $N$ an $S\utens S$-module with involution $\sigma$. Then the orbit Mackey functor
\[
N^{\qdr}=\big(N_{\Ct}\xrightarrow{\id}N_{\Ct}\xrightarrow{\id+\sigma}N^{\Ct}\big)
\]
is a form parameter over $S$, called the \defi{quadratic} form parameter, where the multiplicative action of $S$ on $N_{\Ct}$ is $s\cdot [n]=[(s\otimes s) n]$, for all $s\in S$ and $n\in N^{\Ct}$. The corresponding Poincaré structure on $\Dperf(R)$ is the genuine quadratic structure.

If $R$ is a commutative ring and $M$ is a commutative $R\utens R$-algebra with involution, $S$ is an $R$-algebra,  and $N$ an $M$-module with involution in $R\utens R$-modules, then $N^{\sym}$ and $N^{\qdr}$ are modules over $M^{\sym}$ in $\FPM$. %
The relation for a module from Lemma \reffour{lemma:modulesFPM} for the quadratic case holds since
\[
(n(r)m)(s\cdot [x])=
(r\otimes r)m[s\otimes s x]=[(rs\otimes rs) mx]=(rs)\cdot (m[x])
\]
for $m\in M^{\Ct}, x\in N, r\in R$ and $s\in S$, and similarly for the symmetric case.
\end{example}




\begin{example}%
\label{example:form_parameter_over_Tate_Poincare}%
Let $k$ be a commutative ring. There is a commutative monoid in $\FPMp$
\[
k^{\fr}=\big(k\xrightarrow{\tau}W_{\langle 2\rangle}(k)\xrightarrow{\rho}k\big)
\]
where $W_{\langle 2\rangle}(k)$ is the ring of $2$-truncated $2$-typical Witt vectors of $k$, so that the underlying set of $W_{\langle 2\rangle}(k)$ is $k\times k$. The transfer $\tau$ is the Verschiebung which sends $a$ to $(0,a)$, $\rho$ is the Frobenius that takes $a$ to $a^2+2b$, and $n\colon k\to W_{\langle 2\rangle}(k)$ is the Teichmüller lift that sends $a\in k$ to $(a,0)$. 

We call this form parameter the \defi{Frobenius form parameter}. The corresponding Poincaré structure is the \defi{genuine Frobenius structure} $\QF^{\g\fr}_k$, which coincides with the truncation $\tau_{\leq 0}\QF^{\tate}_k$ of the Tate Poincaré structure 
on $\Dperf(k)$, associated to the module with genuine involution $(k,k,k\to k^{\tC})$ whose structure map is the Tate-valued Frobenius. 
In the case of $k=\ZZ$ the ring $W_{\langle 2\rangle}(\ZZ)$ is isomorphic to the Burside ring and the Frobenius form parameter recovers the Burnside form parameter of Example~\refthree{example:dotto-ogle}.

This construction can also be adapted to the case where $k$ is a commutative ring with involution $\tau$ by replacing $W_{\langle 2\rangle}(k)$ with the ring $W_{\langle 2\rangle}(k,\tau)$ whose underlying set of elements is $k \times k_{\Ct}$ and ring structure defined similarly to the one on $W(k)$ in Example~\reffour{example:mackey}.
\end{example}









Let $k$ be a commutative ring. By Example \reffour{example:fixeptmonoid} the symmetric form parameter $k^{\sym}=(k\xrightarrow{2}k\xrightarrow{\id}k)$ is canonically a monoid in $\FPMp$.


\begin{lemma}
\label{lemma:moduleskg}%
A module over $k^{\sym}$ in $\FPMp$ consists of a form parameter
$
Y=\big(N_{\Ct}\xrightarrow{\tau}P\xrightarrow{\rho}N^{\Ct}\big)
$
over a ring $S$ together with a $k$-algebra structure on $S$ and a $k$-module structure on $P$, such that:
\begin{enumerate}
\item the $S\otimes S$-module structure on $N$ descends to a $S\otimes_k S$-module structure (so that $N$ is in particular a $k$-module and $\sigma$ is $k$-linear),
\item  $\tau$ and $\rho$ are $k$-linear,
\item the action of $S$ on $P$ of the form parameter is $k$-quadratic, in the sense that
\[
(as)\cdot p=a^2(s\cdot p)
 \]
for all $a\in k, s\in S$ and $p\in P$ where $a^2$ on the right acts via the $k$-module structure on $P$.
\end{enumerate}
\end{lemma}
 
\begin{proof}
From the characterisation of modules from Lemma \reffour{lemma:modulesFPM} we see that a $k^{\sym}$-module in $\FPM$ consists of a form parameter $Y$ as above over $S$ where $S$ is a $k$-algebra, $N$ is a $S\otimes_k S$-module, which is a module over the Green functor $k^{\sym}$ and where the action of $S$ on $P$ of the form parameter is $k$-quadratic. Moreover, by Remark \reffour{remark:green} the structure of a module over the Green functor $k^{\sym}$ is the same as a $k$-action on $P$ where $\tau$ and $\rho$ are $k$-linear (for $\tau$ this is because the restriction map of $k^{\sym}$ is surjective) and such that $\tau\rho =2$. We therefore only need to verify this last identity. By combining the formula for the action of the form parameter on a sum with the fact taht this action is $k$-quadratic we have that
 
\[
4(s\cdot p)=(2s)\cdot p=2(s\cdot p)+\tau((s\otimes s)\rho (p)),
\]
 which for $s=1$ gives $\tau(\rho (p))=2p$ for every $p\in P$.
\end{proof}



\begin{remark}
\label{remark:criterion-genuine-symmetric}%
By Lemma \reffour{lemma:modulesFPM}, a $\mathbb{Z}^{\sym}$-module structure on a form parameter $
Y=\big(N_{\Ct}\xrightarrow{\tau}P\xrightarrow{\rho}N^{\Ct}\big)
$
over a ring $S$ exists (and it is in this case unique) if and only if
\[
(as)\cdot p=a^2(s\cdot p)
\]
for all $a\in\mathbb{Z}, s\in S$ and $p\in P$. Moreover, this last property is in fact equivalent to the condition that $\tau\rho=2$. We saw in the proof of Lemma \reffour{lemma:moduleskg} that it is sufficient. Conversely, assume that $\tau\rho=2$. Since $\rho$ is $S$-equivariant we have that
 
\[
2(s\cdot p)=\tau\rho(s\cdot p)=\tau((s\otimes s) \rho (p)),
\]
 and by the property of the action of the form parameter on a sum that
\[
(2s)\cdot p=2(s\cdot p)+\tau((s\otimes s) \rho (p))=4(s\cdot p)=(2)^2(s\cdot p).
\]
By applying this argument inductively on $a$ we can prove that $(as)\cdot p=a^2(s\cdot p)$ for every $a\geq 2$. The formula for $a=0,1$ is obvious. Since
 
\[
0=(as+(-as))\cdot p=(as)\cdot p+(-as)\cdot p+\tau(as\otimes (-as)\rho(p))=(as)\cdot p+(-as)\cdot p-2((as)\cdot p)
\]
 we see that $(-as)\cdot p=(as)\cdot p$ and the formula also holds for negative $a\in\mathbb{Z}$.
\end{remark}

 
 
\begin{lemma}
A module over $k^{\g\fr}$ in $\FPM$ consists of a form parameter
$
Y=\big(N_{\Ct}\xrightarrow{\tau}P\xrightarrow{\rho}N^{\Ct}\big)
$
over a ring $S$ together with a $k$-algebra structure on $S$, such that the $S\otimes S$-module structure on $N$ descends to a $S\otimes_k S$-module structure.
\end{lemma}
 
 
 \begin{proof}
 An action of $W_{\langle 2\rangle}(k)$ on $P$ giving $Y$ the structure of a module is  uniquely determined by the $S$-action of the form parameter. Indeed by setting $s=1$ and $(a,b)=(1,0)$ in the condition
\[
((c,0)(a,b))(s\cdot p)=(cs)\cdot ((a,b)p)
\]
we find that $(a,0)p=a\cdot p$, and from the Frobenius reciprocity condition that
\[
(0,b)p=\tau(b\rho(p)),
\]
for every $a,b\in k$ and $p\in P$. One can directly verify that by defining the action of $W_{\langle 2\rangle}(k)$ on $P$ this way, all the conditions of a module are satisfied.
 \end{proof}
 
 
 
 
\subsection{Form parameters of quadratic and symmetric types}

The next definition captures the two flavours of form parameter originally considered by Bak in \cite{bak-form-parameter}.
 
 \begin{definition}
We call a form parameter $\big(N_{\Ct}\xrightarrow{\tau}P\xrightarrow{\rho}N^{\Ct}\big)$ over a ring $S$ of symmetric type if $\rho$ is injective, and of quadratic type if $\tau$ is surjective. We denote the respective full subcategories of $\FPM$ by $\FPM_\s$ and $\FPM_\qdr$.
 \end{definition}
 
  
\begin{example}\
\begin{enumerate}
\item
Any form parameter of symmetric type is uniquely a module over $\mathbb{Z}^{\sym}$. Indeed, since $\rho$ is injective one can check that the last condition of Lemma \reffour{lemma:moduleskg} holds after applying $\rho$, which is
\[
\rho\tau\rho=(\id+\sigma)\rho=\rho+\sigma\rho=2\rho=\rho 2.
\]
For example the form parameter $M^{\sym}$ associated to an $R$-module with involution is symmetric, and thus a module over $\mathbb{Z}^{\sym}$.
\item
Any form parameter of quadratic type is also uniquely a module over $\mathbb{Z}^{\sym}$. Since $\tau$ is surjective one can choose, for every $p\in P$, an element $n\in N_{\Ct}$ such that $\tau(n)=p$. Thus
\[
\tau\rho(p)=\tau\rho\tau(n)=\tau(n+\sigma(n))=\tau(n)+\tau\sigma(n)=2\tau(n)=2p.
\]
For example the form parameter $M^{\qdr}$ associated to an $R$-module with involution is quadratic, and thus a module over $\mathbb{Z}^{\sym}$. 
\end{enumerate}
\end{example}


We conclude this section with a comparison with the tensor product of form parameters of symmetric type introduced by Bak. Given any module with involution $M$, we write $-M$ for the same module with the involution changed by a sign.


\begin{proposition}%
\label{proposition:symmetric_form_parameters}%
\
\begin{enumerate}
 \item 
The subcategory $\FPM_{\s}$ is a multiplicative Bousfield localisation of $\FPM$, where the left adjoint $s$ of the inclusion is given by
\[
s\big(M_{\Ct}\xrightarrow{\tau}Q\xrightarrow{\rho}M^{\Ct}\big)=\big(M_{\Ct}\xrightarrow{\id+\sigma}Im(\rho)\xrightarrow{}M^{\Ct}\big)
\]
where both form parameters are over the same ring $R$.
In particular $\FPM_{\s}$ inherits a symmetric monoidal structure with unit $\s(\mathbb{A})= \ZZ^{\sym}$. 
\item The action of $\FPM$ onto itself restricts to an action on $\FPM_{\qdr}$, which in turn extends (uniquely) to an action of $\FPM_{\s}$ on $\FPM_{\qdr}$.
\item Tensoring with the form parameter $(-\mathbb{Z})^{\sym}$ gives a functor $(-\mathbb{Z})^{\sym}\boxutens(-)\colon \FPM\to\FPM_{\qdr}$ which moreover identifies with the dual form parameter
\[
 (-\mathbb{Z})^{\sym}\boxutens\fpm=\big((-M)_{\Ct}\xrightarrow{\pr}M/Im(\rho)\xrightarrow{\id-\sigma}(-M)^{\Ct}\big)=\check{\fpm}
\]
from \cite{SchlichtinghigherI} (see also Proposition \refone{proposition:form-parameter-shift}).
\end{enumerate}
\end{proposition}

\begin{proof}
For the first part, it is easy to see that $s$ is left-adjoint to the inclusion. We need to show that tensoring with a form parameter $X=(M_{\Ct}\xrightarrow{\tau}Q\xrightarrow{\rho}M^{\Ct})$ over $R$ preserves $s$-equivalences. These are the morphisms of form parameters $f\colon Y\to Y'$ whose underlying map of rings is an isomorphism, and whose underlying map of modules with involution $g\colon N\to N'$ is an isomorphism and it restricts to a surjection $Im(\rho)\to Im(\rho')$. One can easily see that these are preserved by $X\boxutens (-)$ using that, by the definition of the box product, the image of the restriction map of $X\boxutens Y'$ is generated by the elements of the form $m\otimes n'+\sigma(m)\otimes\sigma(n')$ and $\rho(q)\otimes\rho(p')$ for $m\in M, n'\in N', q\in Q$ and $p'\in P'$. The statement about the unit is clear, since the restriction map $\rho\colon\mathbb{A}(\Ct)\to\mathbb{Z}$ of $\mathbb{A}$ is surjective.



For the second part, let us show that if $X$ is any form parameter and $Y$ is one of quadratic type, then $X\boxutens Y$ is quadratic. By the definition of the transfer of the box product, we need to show that any element of the form $q\otimes p$ for $q\in Q$ and $p\in P$ is equivalent to one of the form $m\otimes n$ for some $m\in M$ and $n\in N$, under the defining equivalence relation of the box product. Since the transfer $\tau$ of $Y$ is assumed to be surjective, there is $n\in N$ with $\tau(n)=p$ and therefore
\[
q\otimes p=q\otimes \tau(n)\sim \rho(q)\otimes n
\]
as we wanted. For the last part, it is sufficient to show that $s$ preserves form parameters of quadratic type, that is that if the transfer $\tau$ of $Y$ is surjective then $\id+\sigma$ surjects onto the image of $\rho$. Indeed, given $p\in P$ we can write $p=\tau(n)$ for some $n\in N$, and therefore
\[
\rho(p)=\rho \tau(n)=n+\sigma(n).
\]


For the final claim, since $(-\mathbb{Z})^{\Ct}=(-\mathbb{Z})_{\Ct}=0$ we find that one of the summands in the definition of the box product $(-\mathbb{Z})^{\sym}\boxutens\fpm$ vanishes, and in the other summand $\mathbb{Z}\otimes M=M$ we have the relation $a\rho(q)\simeq \tau(a)q=0$ for $a\in\mathbb{Z}$ and $q\in Q$. From this one can easily see that $(-\mathbb{Z})^{\sym}\boxutens\fpm\cong \check{\fpm}$.
\end{proof}

\begin{remark}[Comparison to form parameters in the sense of Bak]
Given a ring $R$ and an $(R\otimes_{\ZZ} R)$-module with involution $(M,\sigma)$, 
forms parameters of symmetric type over $R$ extending $(M,\sigma)$ 
are classified by subgroups 
\[
\Lam \subseteq \widehat{\mathrm{H}}^0(\Ct;M) = \mathrm{coker}[\mathrm{N}\colon M_{\Ct} \to M^{\Ct}]
\]
which are stable under the operation $[m] \mapsto [(r\otimes r)(m)]$ for $m \in M^{\Ct}$. Here, the correspondence between these two types of data is given by $\Lambda =  \coker[M_{\Ct}\to Q]$ and $Q = M^{\Ct} \times_{\widehat{\mathrm{H}}^0(\Ct;M)} \Lam$.
 
 Such a triple $(R, M, \Lambda)$ %
 constitutes a form parameter in the sense of Bak (where Bak only considered the case where $R$ is a ring with involution and $M=\pm R$). Under the classification, the reflection $\s\colon \FPM \to \FPM_s$ (and in fact the reflection $\s\circ \tau_{\leq 0}$) are described by sending a module with genuine involution $(M,b\colon N\to M^\tC)$ over $R$ to $(R, M, \im(\pi_0(b))$. Thus, the tensor product in $\FPM_\s$ is described by Bak's tensor product on form parameters,
\[
(R, B, \Lambda)\utens (R', B', \Lambda') = (R\utens R', B\utens B', \im(\Lambda\utens \Lambda'\to \widehat{\mathrm{H}}^0(\Ct, B\utens B')).
\]

We also find that Bak-style form parameter $(R, B, \Lambda)$ classify quadratic-type form parameters just as well, where the duality however needs to be twisted by a sign. The passage from the symmetric-type parameter to the corresponding quadratic-type parameter corresponds to the passage to the ``dual form parameter'' from \S\refone{subsection:discrete-rings}. Under these correspondences, the action $\FPM_\s\times \FPM_\qdr\to \FPM_\qdr$ is again described by Bak's tensor product of form parameters. This explains the tensor product written down by Karoubi \cite{karoubi-periodicity}*{4.3} in our more categorical terms. 
\end{remark}
